* COMMENT latex-header

#+BEGIN_SRC latex
\usepackage{fontspec}
\setmainfont{Linux Libertine O} % libertine definitely better than Times New Roman
% \setmainfont{Times New Roman}
\usepackage[a4paper,width=170mm,top=15mm,bottom=25mm]{geometry}
\usepackage{hyperref} % enable links
\usepackage{ulem} %This is needed to enable underlined emphasis (e.g. _underlined_)
\usepackage{multicol}
\usepackage{graphicx} % enable pictures
\raggedcolumns
\begin{document}
\begin{multicols}{2}
\noindent
\vspace{0.5cm}
ΑΙΤΗΣΗ\\
\noindent
Ιωάννης Ζάννος\\
Καθηγητής Τμήματος Τεχνών Ήχου και Εικόνας\\
Σχολής Μουσικής και Οπτικοακουστικών Τεχνών\\
Ιονίου Πανεπιστημίου

\vspace*{0.5cm}
\noindent
Κέρκυρα, 17/01/2019\\

% \vspace{1cm}
% \noindent
% Κέρκυρα, 17/01/2019\\

\columnbreak
\noindent
ΠΡΟΣ

\vspace{0.5cm}
\noindent
Την Γενική Συνέλευση\\
του Τμήματος Τεχνών Ήχου και Εικόνας\\
Σχολής Μουσικής και Οπτικοακουστικών Τεχνών\\
του Ιονίου Πανεπιστημίου

\vspace{1cm}

#+END_SRC

* COMMENT latex-footer

#+BEGIN_SRC latex
\vspace{1cm}
% \noindent
\parindent=0.5cm
Κέρκυρα,  17. Ιανουαρίου 2019.

\vspace{0.5cm}

\includegraphics[width=3.5cm]{signature}

\vspace{0.5cm}
\parindent=1cm
Ιωάννης Ζάννος

\end{multicols}
\end{document}
#+END_SRC

